\documentclass[11pt]{article}
\usepackage[margin=1in]{geometry}
\usepackage{booktabs}
\usepackage{amsmath}
\usepackage{amssymb}
\usepackage{hyperref}
\usepackage{cite}

\title{Antipsychotic Drugs Selectively Decorrelate Long-Range Interactions in Deep Cortical Layers}
\author{}
\date{}

\begin{document}
\maketitle

\begin{abstract}
Psychosis is characterized by a diminished ability to distinguish externally driven from self-generated activity patterns. Antipsychotic drugs ameliorate psychosis but their circuit-level effects remain unclear. Using widefield calcium imaging of mouse dorsal cortex in a closed-loop/open-loop visual virtual reality paradigm, we quantified cell type-specific activity across six regions of interest (V1, V2am, RSC, M1, A24b, M2). In Tlx3-positive layer 5 intratelencephalic (L5 IT) neurons, closed- and open-loop locomotion onset activity patterns differed markedly, with the similarity measured across all regions being among the lowest of tested cell types. A single intraperitoneal injection of clozapine fundamentally altered L5 IT locomotion-related activity in V1 and produced a reduction in activity correlations across 12 regions (6 per hemisphere) that was significantly stronger for long-range than short-range interactions. Comparing naive C57BL/6 brain-wide GCaMP expression with Tlx3-Cre $\times$ Ai148 L5 IT expression, clozapine effects differed (C57BL/6 vs Tlx3-Cre $\times$ Ai148, short-range: $p < 0.05$; long-range: $p < 10^{-5}$; rank-sum test). Comparing to Cux2-CreERT2 superficial layers (Cux2-CreERT2 $\times$ Ai148 vs Tlx3-Cre $\times$ Ai148, short-range: $p < 0.005$, long-range: $p < 10^{-8}$; rank-sum test). Two additional antipsychotics, aripiprazole and haloperidol, produced analogous decorrelation in L5 IT neurons, while the psychostimulant amphetamine did not. Mismatch responses in V2am were strongly attenuated by antipsychotic drugs while grating responses were less affected. These results support the interpretation that a major functional effect of antipsychotic drugs is a selective alteration of long-range L5-mediated communication.
\end{abstract}

\section{Introduction}
Self-generated movements are coupled to sensory feedback across modalities, and this coupling is believed to support the brain's internal models of the world \cite{jordan1992forward}. Hallucinations and delusions may arise from misconfiguration of these models, producing self-generated activity that is not correctly distinguished from externally driven activity \cite{fletcher2009perceiving,frith2005consciousness}. Locomotion produces broad cortical activation \cite{musall2019singletrial,stringer2019spontaneous} including primary visual cortex \cite{keller2012sensorimotor,saleem2013integration,pakan2016behavioral}. Predictive processing proposes that these movement signals are internal-model predictions of sensory consequences, compared against bottom-up input to yield prediction errors \cite{jiang2021predictive,keller2018predictive}. Prediction error responses have been identified across species \cite{attinger2017visuomotor,audette2021modulation,ayaz2019layerspecific,eliades2008forward,heindorf2018mouse,keller2012sensorimotor,keller2009neural,stanley2007motor,zmarz2016mismatch}; top-down motor inputs carry sensory predictions \cite{leinweber2017sensorimotor,schneider2018movement}. Prediction-error neurons are thought to reside in L2/3 \cite{jordan2020layer,bastos2012canonical}, with internal-representation neurons speculated to lie in deep layers.

Using a closed-loop vs open-loop visuomotor virtual reality paradigm, we asked whether locomotion onsets in mouse dorsal cortex are differentially activated when coupled vs uncoupled from visual flow feedback, and whether antipsychotic drugs modulate these differences in a cell-type-specific way.

\section{Methods}
\subsection{Animals and GCaMP expression}
GCaMP variants were expressed in C57BL/6 (brain-wide via AAV-PHP.eB), Emx1-Cre, or by crossing cell-type-specific Cre driver lines (Cux2-CreERT2, Scnn1a-Cre, Tlx3-Cre, Ntsr1-Cre, PV-Cre, VIP-Cre, SST-Cre) with Ai148 GCaMP6 reporter line \cite{chan2017engineered,daigle2018suite}. A crystal skull cranial window \cite{kim2016longterm} was implanted prior to imaging. Mice were head-fixed on a spherical treadmill and free to locomote \cite{wekselblatt2016largescale}.

\subsection{Experimental conditions}
Mice experienced closed-loop (locomotion controlled forward motion in a virtual corridor) and open-loop (replay of previously self-generated flow) conditions, and relative darkness. Visuomotor mismatch was a 1 s halt of visual flow during closed-loop locomotion. A grating session with drifting gratings covering the entire toroidal screen probed bottom-up visual responses. Six dorsal-cortex ROIs per hemisphere were selected: V1, V2am, RSC, M1, A24b, M2.

\subsection{Drug administration}
A single intraperitoneal injection of clozapine, aripiprazole, haloperidol, or amphetamine, or saline control, was given. Hemodynamic occlusion controls used eGFP-expressing mice. Imaging employed widefield (470 nm LED, sCMOS) and, for confirmation, two-photon calcium imaging of somatic activity in V1.

\subsection{Analyses}
Closed/open-loop similarity was quantified as the correlation coefficient of locomotion-onset activity traces in a $-5$ to $+3$ s window around onset. Inter-regional activity correlations were computed across the 12 ROIs and plotted as a function of bregma-lambda-normalized distance. Short- vs long-range cutoff was set at 0.9$\times$ the bregma-lambda distance ($\approx$3.8 mm). Statistical tests used hierarchical bootstrap and rank-sum tests with family-wise error correction.

\section{Results}
\subsection{L5 IT neurons distinguish closed- and open-loop visuomotor coupling}
In C57BL/6 mice with brain-wide GCaMP expression, closed- and open-loop locomotion onsets produced similar broad activation (an example shows 83 closed-loop and 153 open-loop onsets). In Tlx3-Cre $\times$ Ai148 mice (L5 IT), closed-loop locomotion onsets activated anterior regions but decreased posterior activity (an example shows 88 closed-loop and 83 open-loop onsets). Closed/open-loop similarity across the 12 ROIs was lower in Tlx3 (L5 IT) and Ntsr1 (L6) than in other neuron types tested; family-wise error-corrected significance thresholds are $p < 0.05/9$, $p < 0.01/9$, and $p < 0.001/9$ for comparisons against Tlx3. Correlation between calcium activity and locomotion speed in C57BL/6 was comparable across closed-loop, open-loop, and dark; in Tlx3 mice, closed-loop correlation was negative in posterior cortex and positive anteriorly, open-loop correlation was positive throughout, and darkness mirrored open-loop posteriorly and closed-loop anteriorly. Sessions counts: 308 closed-loop, 316 open-loop, 68 dark in C57BL/6; 420 closed-loop, 394 open-loop, 194 dark in Tlx3-Cre $\times$ Ai148.

\begin{table}[h]
\centering
\caption{Recording and stimulation session counts across preparations.}
\begin{tabular}{lll}
\toprule
Genotype & Closed/open/dark (5-min sessions) & Notes \\
\midrule
C57BL/6 (brain-wide, $n=6$) & 308 / 316 / 68 & GCaMP via AAV-PHP.eB \\
Tlx3-Cre $\times$ Ai148 (L5 IT, $n=15$) & 420 / 394 / 194 & GCaMP6 in L5 IT \\
\bottomrule
\end{tabular}
\end{table}

\subsection{Visuomotor prediction errors originate in visual cortex and drive L5 IT differentially}
In C57BL/6 mice, visuomotor mismatch (230 onsets) and drifting grating onsets (86 onsets) produced activity increases first and strongest in visual regions. In Tlx3-Cre $\times$ Ai148 (L5 IT), mismatch (292 onsets) elicited an increase in V1 activity while grating onsets (171 onsets) elicited a decrease in V1 activity. This pattern is consistent with either negative prediction error encoding or internal-representation-like encoding.

\subsection{Clozapine alters L5 IT locomotion-related activity in V1}
In C57BL/6 mice, clozapine had minimal effect on closed- and open-loop locomotion onset responses, with a small increase in open-loop visual flow responses. In Tlx3-Cre $\times$ Ai148 (L5 IT), clozapine caused massive increases in V1 locomotion-related activity for both closed- and open-loop conditions; open-loop visual flow onset responses remained unchanged. The effect on locomotion vs activity correlation was opposite in the two genotypes: increased in C57BL/6 but decreased in Tlx3. Two-photon imaging confirmed the widefield findings at the somatic level.

\subsection{Antipsychotic-induced decorrelation is strongest in L5 IT and for long-range pairs}
Correlations among 12 ROIs were computed in 4 C57BL/6 mice and 5 Tlx3-Cre $\times$ Ai148 mice before and after clozapine. Data were split into short- and long-range portions at 0.9$\times$ the bregma-lambda distance ($\approx$3.8 mm). In C57BL/6, reductions in short- (122 region pairs) and long-range (142 region pairs) correlations were not significant. In Tlx3-Cre $\times$ Ai148, both short- (182 pairs) and long-range (148 pairs) correlations were significantly reduced, with stronger effects on long-range pairs ($p < 0.001$). Comparing genotypes: C57BL/6 vs Tlx3-Cre $\times$ Ai148, short-range $p < 0.05$; long-range $p < 10^{-5}$; rank-sum test. Comparing superficial (Cux2-CreERT2 $\times$ Ai148) vs deep (Tlx3-Cre $\times$ Ai148): short-range $p < 0.005$; long-range $p < 10^{-8}$; rank-sum test. Saline controls showed no change. Hemodynamic-occlusion controls in eGFP-expressing Tlx3-Cre mice could not account for the decrease and produced a small increase if anything. Two-photon pairwise somatic activity correlations did not change.

\begin{table}[h]
\centering
\caption{Clozapine effects on short/long-range activity correlations by genotype.}
\begin{tabular}{llll}
\toprule
Genotype & Short-range pairs & Long-range pairs & Effect vs naive \\
\midrule
C57BL/6 brain-wide ($n=4$) & 122 & 142 & n.s. \\
Tlx3-Cre $\times$ Ai148 (L5 IT, $n=5$) & 182 & 148 & significant, $p < 0.001$ long-range \\
C57BL/6 vs Tlx3 (rank-sum) & $p < 0.05$ & $p < 10^{-5}$ & \\
Cux2 vs Tlx3 (rank-sum) & $p < 0.005$ & $p < 10^{-8}$ & \\
\bottomrule
\end{tabular}
\end{table}

\subsection{Aripiprazole and haloperidol replicate decorrelation; amphetamine does not}
In 3 Tlx3-Cre $\times$ Ai148 mice, aripiprazole produced significant decorrelation across 114 short-range and 84 long-range pairs ($*$: $p < 0.05$; $***$: $p < 0.001$). Haloperidol in 3 additional Tlx3-Cre $\times$ Ai148 mice produced analogous effects across 114 and 84 pairs ($*$: $p < 0.05$; long-range n.s. in some comparisons). Amphetamine in 3 Tlx3-Cre $\times$ Ai148 mice produced no significant change in 114 short-range and 84 long-range pairs (n.s.).

\begin{table}[h]
\centering
\caption{Antipsychotic drug effects on L5 IT correlations.}
\begin{tabular}{lllll}
\toprule
Drug & Mice & Short-range pairs & Long-range pairs & Effect \\
\midrule
Clozapine & 5 & 182 & 148 & significant short + long \\
Aripiprazole & 3 & 114 & 84 & significant ($p < 0.05$, $p < 0.001$) \\
Haloperidol & 3 & 114 & 84 & short $p < 0.05$; long n.s. \\
Amphetamine & 3 & 114 & 84 & n.s. \\
\bottomrule
\end{tabular}
\end{table}

\subsection{Antipsychotic drugs attenuate prediction-error spread from V1 to V2am}
Averaged across clozapine (5 mice), aripiprazole (3 mice), and haloperidol (3 mice), visuomotor mismatch responses in Tlx3-Cre $\times$ Ai148 were partially reduced in V1 and almost absent in V2am. Drifting grating responses in V2am were also reduced, but overall less affected than mismatch responses.

\section{Discussion}
We show that antipsychotic drugs primarily alter cortical function by decorrelating long-range activity in L5 IT neurons, while leaving superficial layers and brain-wide averages comparatively intact. Closed- vs open-loop visuomotor contrasts in Tlx3-positive L5 IT neurons are consistent with an internal-representation role for these neurons \cite{leinweber2017sensorimotor,bastos2012canonical,keller2018predictive}, and clozapine's massive enhancement of V1 locomotion-related activity supports the idea that antipsychotics reduce the precision of internal predictions.

Because L5 IT neurons are a primary source of long-range cortico-cortical communication \cite{chen2019empty,harris2019hierarchical,leinweber2017sensorimotor}, the stronger reduction in long-range pairs is consistent with a selective attenuation of distal L5-mediated signaling. The replication with aripiprazole and haloperidol---but absence under amphetamine---indicates this effect is a shared signature of antipsychotics rather than a generic sedative or arousal change. The preferential attenuation of mismatch responses over grating responses suggests the drugs primarily affect the spread of top-down-mediated prediction errors, consistent with L5-mediated lateral and long-range communication in cortex \cite{jordan2020layer}.

Limitations include the approximation of axonal path length by top-down linear distance, the inability to separate direct L5-L5 communication from subcortical loops, and the single-dose design. Future work will test whether repeated antipsychotic administration produces cumulative or plastic changes in L5 decorrelation, and whether optogenetic disruption of L5 IT long-range projections phenocopies the antipsychotic signature.

\bibliographystyle{plain}
\bibliography{references}

\end{document}
